\documentclass[11pt]{amsart}

\usepackage[margin=1.15in]{geometry}
\usepackage{amsmath,amssymb,amsthm,mathtools}
\usepackage{enumitem}
\usepackage{microtype}
\usepackage[colorlinks=true,linkcolor=blue,citecolor=blue,urlcolor=blue]{hyperref}
\usepackage[nameinlink,capitalize]{cleveref}

\theoremstyle{plain}
\newtheorem{theorem}{Theorem}[section]
\newtheorem{proposition}[theorem]{Proposition}
\newtheorem{lemma}[theorem]{Lemma}
\newtheorem{corollary}[theorem]{Corollary}

\theoremstyle{definition}
\newtheorem{definition}[theorem]{Definition}
\newtheorem{assumption}[theorem]{Assumption}

\theoremstyle{remark}
\newtheorem{remark}[theorem]{Remark}

\numberwithin{equation}{section}

\DeclareMathOperator{\diver}{div}

\title[Lyapunov branch exclusions for ancient Navier--Stokes flows]
{Lyapunov Branch Exclusions for Ancient Navier--Stokes Flows}
\author{Guillem Duran Ballester}
\date{}

\subjclass[2020]{35Q30, 35B40, 35B45, 76D05}
\keywords{Navier--Stokes equations, ancient solutions, axisymmetric flows, Lyapunov functionals, Burgers vortices}

\begin{document}

\begin{abstract}
We isolate the parts of a Lyapunov-based ancient-solution analysis for the
three-dimensional Navier--Stokes equations which presently give rigorous
branch exclusions.  For the unforced equations, bounded ancient axisymmetric
subclasses are eliminated by known Liouville theorems: no swirl, pointwise
scale-invariant control, finite \(L^p\) control of the swirl variable
\(\Gamma=ru_\theta\), and axial periodicity with bounded \(\Gamma\).  A
separate perturbative weighted-vorticity branch is excluded by a semiflow
Lyapunov functional built from Gallay--Wayne's small-data decay theorem in weighted
vorticity spaces.  For the strained Burgers-vortex model, the
Gallay--Maekawa stability theorem yields a profile-relative Lyapunov functional
and an ancient rigidity theorem near the Burgers-vortex family.  Finally, in
the open bounded-swirl axisymmetric case, we prove the exact scalar
\(\Gamma\)-defect identity and the anisotropic shell-error decomposition.  In
the whole-space smooth bounded setting, the clean global identity is centered
at the only constant compatible with regularity on the axis, namely
\(\Gamma=0\).  This identity does not close the open branch by itself; it gives
a precise coercive shell-control criterion under which the branch would reduce
to vanishing \(\Gamma\).
\end{abstract}

\maketitle

\section{Introduction}

We consider ancient solutions of the incompressible Navier--Stokes equations
on \(\mathbb R^3\),
\begin{equation}\label{eq:ns}
   \partial_t u+(u\cdot\nabla)u+\nabla p=\Delta u,\qquad
   \nabla\cdot u=0,
   \qquad t<0.
\end{equation}
The purpose of this paper is not to claim a new Liouville theorem for all
bounded ancient solutions.  Rather, it records exactly which branches are
closed by existing Liouville theorems or by a genuine Lyapunov construction,
and it writes the remaining bounded-swirl obstruction as an explicit
shell-transport problem.

There are three distinct mechanisms below.

First, several axisymmetric branches of the unforced equations are already
closed by known Liouville theorems.  If a bounded ancient mild solution is
axisymmetric without swirl, Koch--Nadirashvili--Seregin--Sverak imply that it
is Galilean-trivial.  The same conclusion holds under the pointwise
scale-invariant bound \(|u|\le C/r\).  Two swirl-bearing subclasses are also
closed: Lei--Zhang--Zhao prove triviality when
\(\Gamma=ru_\theta\in L_t^\infty L_x^p\), \(1\le p<\infty\), and
Lei--Ren--Zhang prove triviality for \(z\)-periodic bounded ancient
axisymmetric solutions with bounded \(\Gamma\).

Second, a perturbative weighted-vorticity branch is closed by an actual
Lyapunov functional.  Gallay--Wayne's small-data theory for the rescaled
three-dimensional vorticity equation gives exponential decay in weighted
\(L^2(m)\) spaces.  From this decay one defines
\[
   \mathcal L(w_0)=\sup_{\tau\ge0}e^{2\tau}\|\Phi_\tau w_0\|_{L^2(m)}^2 .
\]
This functional is coercive on the small ball and contracts along the forward
semiflow.  A bounded ancient trajectory contained in that ball is therefore
forced to vanish by letting the initial time tend to \(-\infty\).

Third, the same Lyapunov construction closes a model structured-vortex branch
for the strained Burgers-vortex equation.  The neutral circulation parameter is
removed by horizontal integration of the vertical vorticity; Gallay--Maekawa's
stability theorem then provides exponential decay on the fixed-circulation
subspace.  The associated profile-relative functional excludes all ancient
solutions which remain in the small weighted tube around the Burgers-vortex
family.  This is a theorem for the strained model.  Importing it into an
unforced Navier--Stokes blow-up analysis requires an additional analytic
reduction to that model.

The last part treats the open bounded-swirl axisymmetric branch.  The correct
first Lyapunov variable is not the full velocity defect, but the scalar swirl
quantity
\[
   \Gamma=ru_\theta .
\]
For \(W=\Gamma\) and for any nonnegative axisymmetric test weight \(\phi\), we
prove
\begin{equation}\label{eq:intro-gamma}
   \frac12\frac{d}{dt}\int_{\mathbb R^3}W^2\phi\,dx
   +\int_{\mathbb R^3}|\nabla W|^2\phi\,dx
   =
   \frac12\int_{\mathbb R^3}W^2
   \left(
      \partial_t\phi+\Delta\phi+b\cdot\nabla\phi+\frac2r\partial_r\phi
   \right)\,dx,
\end{equation}
where \(b=u_r e_r+u_z e_z\).  For \(W=\Gamma-c\), the same formula has an
additional axis boundary term unless \(c=0\) or the test weight avoids the
axis.  With
\(\phi(r,z,t)=\lambda_R(r)\eta_Z(z-\zeta(t))\), this becomes an exact
anisotropic shell formula whose only error channels are radial shell transport,
axial drift mismatch, and cutoff curvature.  The radial diffusion operator is
\(\partial_r^2+3r^{-1}\partial_r\), so the adapted radial barrier is harmonic
for the four-dimensional radial Laplacian, not an isotropic Gaussian.

The identity \eqref{eq:intro-gamma} is rigorous and useful, but it does not
prove the full bounded-\(\Gamma\) Liouville conjecture.  We therefore state the
precise coercive shell-control criterion that would force \(\Gamma\) to vanish
in expanding windows.  The open analytic point is to verify that
criterion for every bounded ancient axisymmetric solution with bounded
\(\Gamma\) and no decay or periodicity.

\subsection*{Relation with the literature}

Bounded ancient solutions enter the Navier--Stokes regularity theory through
blow-up and compactness arguments.  The backward-uniqueness method of
Escauriaza--Seregin--Sverak \cite{ESS2003} proves regularity for
\(L^\infty_tL^3_x\) solutions and is one of the standard mechanisms by which
critical control rules out singularity formation.  In the axisymmetric setting,
Koch--Nadirashvili--Seregin--Sverak \cite{KNSS2009} developed Liouville
theorems for bounded ancient solutions and connected them to lower bounds for
axisymmetric singularity formation.  The related quantitative lower-bound
program for axisymmetric solutions with swirl was advanced by
Chen--Strain--Tsai--Yau \cite{CSTY2008,CSTY2009}, while
Seregin--Sverak \cite{SereginSverak2009} treated Type I local
axisymmetric singularities.

The finite-\(L^p\) swirl branch used here is the bounded ancient mild theorem
of Lei--Zhang--Zhao \cite{LZZ}.  It should be distinguished from the earlier
Lei--Zhang theorem \cite{LeiZhang2011}, which proves regularity and Liouville
consequences under scale-invariant control of the meridional drift
\(b=u_re_r+u_ze_z\), for instance \(b\in L^\infty_t\mathrm{BMO}^{-1}_x\).
The periodic bounded-swirl branch is the theorem of Lei--Ren--Zhang
\cite{LRZ} on \(\mathbb R^2\times\mathbb T^1\).  These results are used below
only in the precise forms stated in \Cref{thm:knss-input,thm:knss-pointwise-input,thm:lzz-input,thm:lrz-input};
the paper does not interpolate between them or assert an unstated Liouville
theorem for the whole nonperiodic bounded-\(\Gamma\) class.

The Lyapunov arguments in the perturbative model sections rely on stability
theory rather than on axisymmetric Liouville theorems.  Gallay--Wayne
\cite{GallayWayne2002} established decay for the rescaled vorticity equation
in weighted spaces under smallness and algebraic spatial decay.  For Burgers
vortices, linear and low-Reynolds-number stability results such as
Schmid--Rossi \cite{SchmidRossi2004} and Gallay--Wayne \cite{GallayWayne2006}
precede the nonlinear three-dimensional stability theorem of Gallay--Maekawa
\cite{GallayMaekawa}, which is the stability input used in
\Cref{thm:gm}.  The Burgers-vortex conclusions in this paper are therefore
model conclusions unless an independent reduction from an unforced
Navier--Stokes regime to the strained equation has been proved.

\section{Preliminaries}

\subsection{Ancient solutions and Galilean equivalence}

\begin{definition}\label{def:ancient}
A velocity-pressure pair \((u,p)\) is a bounded ancient mild solution of
\eqref{eq:ns} if it solves \eqref{eq:ns} on
\(\mathbb R^3\times(-\infty,0)\) in the mild sense on every finite time
interval, namely for every \(-\infty<s<t<0\),
\begin{equation}\label{eq:mild-ns}
   u(t)=e^{(t-s)\Delta}u(s)
   -\int_s^t e^{(t-\sigma)\Delta}\mathbb P\nabla\cdot
      (u(\sigma)\otimes u(\sigma))\,d\sigma
\end{equation}
in distributions, where \(\mathbb P\) is the Leray projector, and
\[
   \|u\|_{L^\infty(\mathbb R^3\times(-\infty,0))}<\infty .
\]
The pressure is understood modulo functions of time.
\end{definition}

\begin{remark}[Regularity used below]\label{rem:regularity}
The branch exclusions in \Cref{thm:no-swirl,thm:pointwise,thm:lp-swirl,thm:periodic}
use the cited Liouville theorems in their stated bounded ancient mild or weak
classes.  The differential identities in \Cref{lem:gamma-equation,prop:gamma-identity,prop:shell}
are proved for smooth axisymmetric solutions.  When they are invoked for
bounded ancient mild axisymmetric solutions, the required additional input is
that the solution is smooth on \(\mathbb R^3\times(-\infty,0)\) and that
\(\Gamma\), \(\nabla\Gamma\), and the displayed weighted integrals are finite
for the chosen compactly supported weights.  These are included explicitly in
the shell-control assumption below, rather than inferred silently.
\end{remark}

\begin{lemma}[Galilean reduction]\label{lem:galilean}
Let \((u,p)\) solve \eqref{eq:ns}.  For every constant vector \(c\in\mathbb R^3\),
\[
   \widetilde u(x,t)=u(x+ct,t)-c,\qquad
   \widetilde p(x,t)=p(x+ct,t)
\]
also solves \eqref{eq:ns}.  In particular, every constant velocity field is
Galilean-equivalent to zero.
\end{lemma}

\begin{proof}
All derivatives below are evaluated at the point \((x+ct,t)\).  By the chain
rule,
\[
   \partial_t\widetilde u
   =
   (\partial_tu+c\cdot\nabla u)(x+ct,t),
   \qquad
   (\widetilde u\cdot\nabla)\widetilde u
   =
   ((u-c)\cdot\nabla u)(x+ct,t).
\]
Adding these two identities gives
\[
   \partial_t\widetilde u+(\widetilde u\cdot\nabla)\widetilde u
   =
   (\partial_tu+(u\cdot\nabla)u)(x+ct,t).
\]
Also
\[
   \Delta\widetilde u(x,t)=(\Delta u)(x+ct,t),
   \qquad
   \nabla\widetilde p(x,t)=(\nabla p)(x+ct,t),
\]
and
\[
   \nabla\cdot\widetilde u(x,t)
   =
   (\nabla\cdot u)(x+ct,t).
\]
Substituting the equation for \((u,p)\) at \((x+ct,t)\) therefore gives
\[
   \partial_t\widetilde u+(\widetilde u\cdot\nabla)\widetilde u
      +\nabla\widetilde p
   =
   \Delta\widetilde u,
   \qquad
   \nabla\cdot\widetilde u=0 .
\]
If \(u\equiv c\), then \(\widetilde u(x,t)=c-c=0\) for every \((x,t)\).
\end{proof}

\subsection{Axisymmetric notation}

Write cylindrical coordinates as
\[
   x=(r\cos\theta,r\sin\theta,z),\qquad
   u=u_r e_r+u_\theta e_\theta+u_z e_z .
\]
The solution is axisymmetric if \(u_r,u_\theta,u_z\) are independent of
\(\theta\).  It is swirl-free if \(u_\theta\equiv0\).  The scale-invariant
swirl variable is
\[
   \Gamma=r u_\theta .
\]
For axisymmetric solutions the meridional drift is
\[
   b=u_r e_r+u_z e_z .
\]

\section{Exact axisymmetric Liouville branch exclusions}

The results in this section are Liouville imports for the actual unforced
Navier--Stokes equations.  They are stronger than model stability statements:
no reduction to an auxiliary equation is involved.

\subsection{External Liouville inputs}

The following results are used as external theorems.  They are stated here in
the precise forms needed in the sequel.

\begin{theorem}[Koch--Nadirashvili--Seregin--Sverak, no swirl]\label{thm:knss-input}
Let \(u\) be a bounded ancient weak axisymmetric solution of \eqref{eq:ns} on
\(\mathbb R^3\times(-\infty,0)\).  If \(u_\theta\equiv0\), then
\[
   u(x,t)=b(t)e_z
\]
for some bounded scalar function \(b(t)\).
\end{theorem}

\begin{proof}
This is the no-swirl Liouville theorem of
Koch--Nadirashvili--Seregin--Sverak \cite{KNSS2009}, stated only in the form
used below.
\end{proof}

\begin{theorem}[Koch--Nadirashvili--Seregin--Sverak, pointwise branch]
\label{thm:knss-pointwise-input}
Let \(u\) be a bounded ancient weak axisymmetric solution of \eqref{eq:ns} on
\(\mathbb R^3\times(-\infty,0)\).  If
\[
   |u(x,t)|\le \frac{C}{r}
   \qquad (r>0,\ t<0),
\]
then \(u\equiv0\).
\end{theorem}

\begin{proof}
This is Theorem 5.3 of \cite{KNSS2009}.
\end{proof}

\begin{theorem}[Lei--Zhang--Zhao]\label{thm:lzz-input}
Let \(u\) be a bounded ancient mild axisymmetric solution of \eqref{eq:ns} on
\(\mathbb R^3\times(-\infty,0)\).  If, for some \(1\le p<\infty\),
\[
   \Gamma=ru_\theta\in L_t^\infty L_x^p
      \bigl(\mathbb R^3\times(-\infty,0)\bigr),
\]
then \(u\) is a constant vector field.
\end{theorem}

\begin{proof}
This is the ancient axisymmetric \(L^p\)-swirl Liouville theorem of
Lei--Zhang--Zhao \cite{LZZ}, stated in the form used here.
\end{proof}

\begin{theorem}[Lei--Ren--Zhang]\label{thm:lrz-input}
Let \(u\) be a bounded ancient mild axisymmetric solution of \eqref{eq:ns} on
\(\mathbb R^2\times\mathbb T^1\).  If \(\Gamma=ru_\theta\) is bounded, then
\[
   u(x,t)=ce_z
\]
for some constant \(c\in\mathbb R\).
\end{theorem}

\begin{proof}
This is the periodic bounded-swirl Liouville theorem of
Lei--Ren--Zhang \cite{LRZ}.
\end{proof}

\subsection{Consequences for branch exclusions}

\begin{theorem}[No-swirl ancient branch]\label{thm:no-swirl}
Let \(u\) be a bounded ancient mild solution of \eqref{eq:ns}.  If \(u\) is
axisymmetric and \(u_\theta\equiv0\), then \(u\) is Galilean-trivial.  More
precisely, after rotation of the notation around the symmetry axis,
\[
   u(x,t)=c e_z
   \qquad\text{for some }c\in\mathbb R .
\]
\end{theorem}

\begin{proof}
\Cref{thm:knss-input} gives \(u(x,t)=b(t)e_z\).  This means that \(u\) is
independent of \(x\) and points along the symmetry axis.  Hence
\(\Delta u=0\), and \((u\cdot\nabla)u=0\).  Substitution into \eqref{eq:ns}
gives
\[
   b'(t)e_z+\nabla p(x,t)=0 .
\]
The preceding identity alone allows the pressure to absorb \(b'(t)e_z\).
The mild formulation fixes the velocity more strongly.  For any
\(-\infty<s<t<0\), \eqref{eq:mild-ns} gives
\[
   u(t)
   =
   e^{(t-s)\Delta}u(s)
   -
   \int_s^t e^{(t-\sigma)\Delta}\mathbb P\nabla\cdot
      (u(\sigma)\otimes u(\sigma))\,d\sigma .
\]
Here \(u(s)=b(s)e_z\) is a constant vector field, so
\[
   e^{(t-s)\Delta}u(s)=u(s)=b(s)e_z .
\]
Moreover \(u(\sigma)\otimes u(\sigma)\) is a constant matrix in \(x\), so
\(\nabla\cdot(u(\sigma)\otimes u(\sigma))=0\) distributionally.  The integral
term is therefore zero.  Thus \(b(t)e_z=b(s)e_z\) for every \(s<t<0\), and
\(b\) is constant.  The resulting constant velocity field is removed by
\Cref{lem:galilean}.
\end{proof}

\begin{theorem}[Pointwise scale-invariant axisymmetric branch]\label{thm:pointwise}
Let \(u\) be a bounded ancient mild axisymmetric solution of \eqref{eq:ns}.
If
\[
   |u(x,t)|\le \frac{C}{r}
   \qquad (r>0,\ t<0),
\]
then \(u\equiv0\).
\end{theorem}

\begin{proof}
The hypothesis makes \(u\) a bounded ancient axisymmetric weak solution
satisfying the pointwise scale-invariant condition in
\Cref{thm:knss-pointwise-input}.  A mild solution is distributional, and
boundedness gives the local integrability required in the weak formulation.
Applying \Cref{thm:knss-pointwise-input} gives \(u\equiv0\).
\end{proof}

\begin{theorem}[\(L^p\)-controlled swirl branch]\label{thm:lp-swirl}
Let \(u\) be a bounded ancient mild axisymmetric solution of \eqref{eq:ns}.
Assume that for some \(1\le p<\infty\),
\[
   \Gamma=ru_\theta\in L_t^\infty L_x^p
      \bigl(\mathbb R^3\times(-\infty,0)\bigr).
\]
Then \(u\) is Galilean-trivial.
\end{theorem}

\begin{proof}
If \(u_\theta\equiv0\), then \(u\) lies in the no-swirl class and
\Cref{thm:no-swirl} gives Galilean triviality.  If \(u_\theta\not\equiv0\),
then the displayed assumption is the finite-\(L^p\) swirl hypothesis in
\Cref{thm:lzz-input}.  Therefore \(u\) is a constant vector field.  By
\Cref{lem:galilean}, every constant vector field is equivalent to zero under a
Galilean transform.
\end{proof}

\begin{theorem}[Periodic bounded-swirl branch]\label{thm:periodic}
Let \(u\) be a bounded ancient mild axisymmetric solution of \eqref{eq:ns}.
Assume that \(u\) is periodic in \(z\) and
\[
   \Gamma=ru_\theta\in L^\infty
      \bigl(\mathbb R^3\times(-\infty,0)\bigr).
\]
Then \(u(x,t)=ce_z\) for some \(c\in\mathbb R\).  Hence \(u\) is
Galilean-trivial.
\end{theorem}

\begin{proof}
Let \(Z_0>0\) be an axial period.  Since \(u(r,\theta,z+Z_0,t)=u(r,\theta,z,t)\),
the solution descends to the quotient
\(\mathbb R^2\times(\mathbb R/Z_0\mathbb Z)\).  The theorem of
Lei--Ren--Zhang is invariant under the Navier--Stokes scaling
\[
   u_\lambda(x,t)=\lambda u(\lambda x,\lambda^2t),
   \qquad
   p_\lambda(x,t)=\lambda^2p(\lambda x,\lambda^2t).
\]
Choosing \(\lambda=Z_0/(2\pi)\) normalizes the axial period to \(2\pi\), so the
quotient is \(\mathbb R^2\times\mathbb T^1\).  Boundedness of \(u\) is
preserved up to the constant factor \(\lambda\), and
\[
   r(u_\lambda)_\theta(r,z,t)
   =
   r\lambda u_\theta(\lambda r,\lambda z,\lambda^2t)
   =
   (\lambda r)u_\theta(\lambda r,\lambda z,\lambda^2t),
\]
so boundedness of \(\Gamma=ru_\theta\) is preserved with the same
\(L^\infty\) bound under this scaling.
\Cref{thm:lrz-input} then gives \(u(x,t)=ce_z\) for some constant \(c\).
\Cref{lem:galilean} removes this constant axial drift.
\end{proof}

\begin{corollary}[Closed unforced axisymmetric branches]\label{cor:axis-closed}
Let \(\mathcal A\) be any family of bounded ancient mild axisymmetric solutions
of \eqref{eq:ns}.  Suppose every \(u\in\mathcal A\) satisfies at least one of
the following conditions:
\begin{enumerate}[label=\textup{(\roman*)}]
\item \(u_\theta\equiv0\);
\item \(|u(x,t)|\le C/r\);
\item \(ru_\theta\in L_t^\infty L_x^p\) for some \(1\le p<\infty\);
\item \(u\) is periodic in \(z\) and \(ru_\theta\in L^\infty\).
\end{enumerate}
Then every element of \(\mathcal A\) is Galilean-trivial.  If the class is
studied in a gauge where constant drifts are fixed to zero, then
\(\mathcal A=\{0\}\).
\end{corollary}

\begin{proof}
Fix \(u\in\mathcal A\).  By hypothesis, \(u\) satisfies at least one of the
four listed alternatives.  In case (i), \Cref{thm:no-swirl} applies.  In case
(ii), \Cref{thm:pointwise} applies.  In case (iii), \Cref{thm:lp-swirl}
applies.  In case (iv), \Cref{thm:periodic} applies.  Each conclusion is
Galilean triviality, and the extra gauge assumption converts Galilean
triviality into the equality \(u=0\).
\end{proof}

\section{A perturbative weighted-vorticity Lyapunov exclusion}

We next record a genuine Lyapunov exclusion for a perturbative weighted branch.
This section is independent of axisymmetry.

\subsection{Rescaled vorticity equation}

Let \(w=\nabla\times v\) be the vorticity in forward self-similar variables.
The Gallay--Wayne rescaled vorticity equation is
\begin{equation}\label{eq:gallay-vorticity}
   \partial_\tau w
   =
   \Lambda w-(v\cdot\nabla)w+(w\cdot\nabla)v,
   \qquad \nabla\cdot w=0,
\end{equation}
where
\begin{equation}\label{eq:lambda}
   \Lambda w=\Delta w+\frac12 \xi\cdot\nabla w+w,
\end{equation}
and \(v\) is recovered from \(w\) by the three-dimensional Biot--Savart law.
For \(m\ge0\), set
\[
   L^2_\sigma(m)
   =
   \left\{
      f\in L^2(\mathbb R^3;\mathbb R^3):
      \nabla\cdot f=0,\quad
      \|f\|_m<\infty
   \right\},
\]
where
\[
   \|f\|_m^2
   =
   \int_{\mathbb R^3}(1+|\xi|)^{2m}|f(\xi)|^2\,d\xi .
\]

\begin{theorem}[Gallay--Wayne small-data decay]\label{thm:gallay}
Let \(0<\mu\le1\) and \(m>2\mu+\frac12\).  There exist constants
\(\rho_m>0\) and \(C_m\ge1\) such that if
\(\|w_0\|_m\le\rho_m\), then \eqref{eq:gallay-vorticity} has a unique global
mild solution \(w(\tau)=\Phi_\tau w_0\) in
\[
   C([0,\infty);L^2_\sigma(m))
\]
and
\[
   \|\Phi_\tau w_0\|_m
   \le
   C_m e^{-\mu\tau}\|w_0\|_m,\qquad \tau\ge0.
\]
In particular, if \(m>\frac72\), one may use \(\mu=1\).
\end{theorem}

\begin{proof}
This is the Gallay--Wayne weighted small-data theorem for the rescaled
three-dimensional vorticity equation \cite{GallayWayne2002}.  It is used as
an external input.  The only conclusions used below are existence, uniqueness,
the semiflow property induced by uniqueness, and the displayed exponential
decay estimate.
\end{proof}

\begin{definition}[Perturbative weighted ancient branch]\label{def:weighted-branch}
Fix \(m>\frac72\), and let \(\rho_m\) be the small-data radius in
\Cref{thm:gallay} with \(\mu=1\).  The perturbative weighted ancient branch is
the set of all
\[
   w\in C(\mathbb R;L^2_\sigma(m))
\]
which solve \eqref{eq:gallay-vorticity} for all \(\tau\in\mathbb R\) and
satisfy
\[
   \sup_{\tau\in\mathbb R}\|w(\tau)\|_m\le\rho_m .
\]
Solving for all \(\tau\in\mathbb R\) means that for every \(s\in\mathbb R\),
the shifted function \(w_s(\theta)=w(s+\theta)\), \(\theta\ge0\), is the
Gallay--Wayne mild solution with initial datum \(w(s)\).
\end{definition}

\begin{proposition}[Semiflow Lyapunov functional]\label{prop:weighted-lyap}
Fix \(m>\frac72\).  For \(\|w_0\|_m\le\rho_m\), define
\begin{equation}\label{eq:weighted-lyap}
   \mathcal L_m(w_0)
   =
   \sup_{\tau\ge0}e^{2\tau}\|\Phi_\tau w_0\|_m^2 .
\end{equation}
Then \(\mathcal L_m\) is well-defined and satisfies
\begin{equation}\label{eq:weighted-coercive}
   \|w_0\|_m^2
   \le
   \mathcal L_m(w_0)
   \le
   C_m^2\|w_0\|_m^2 .
\end{equation}
Moreover, for every \(s\ge0\),
\begin{equation}\label{eq:weighted-contract}
   \mathcal L_m(\Phi_s w_0)
   \le
   e^{-2s}\mathcal L_m(w_0).
\end{equation}
\end{proposition}

\begin{proof}
By \Cref{thm:gallay} with \(\mu=1\),
\[
   e^{2\tau}\|\Phi_\tau w_0\|_m^2
   \le
   C_m^2\|w_0\|_m^2
   \qquad(\tau\ge0),
\]
so the supremum in \eqref{eq:weighted-lyap} is finite and the upper bound in
\eqref{eq:weighted-coercive} follows.  The lower bound follows by choosing
\(\tau=0\), since \(\Phi_0w_0=w_0\).

For the contraction estimate, first observe that the solution starting from
\(\Phi_sw_0\) at time \(0\) is the same as the original solution starting from
\(w_0\) and observed after the additional time \(s\).  This is the
semiflow property following from uniqueness:
\[
   \Phi_\tau(\Phi_sw_0)=\Phi_{\tau+s}w_0,\qquad \tau,s\ge0.
\]
Therefore
\[
   \mathcal L_m(\Phi_s w_0)
   =
   \sup_{\tau\ge0}e^{2\tau}\|\Phi_{\tau+s}w_0\|_m^2 .
\]
Writing \(\sigma=\tau+s\), the range \(\tau\ge0\) becomes \(\sigma\ge s\), and
also \(e^{2\tau}=e^{2(\sigma-s)}=e^{-2s}e^{2\sigma}\).  Hence
\[
   \mathcal L_m(\Phi_s w_0)
   =
   e^{-2s}\sup_{\sigma\ge s}e^{2\sigma}\|\Phi_\sigma w_0\|_m^2
   \le
   e^{-2s}\mathcal L_m(w_0).
\]
\end{proof}

\begin{theorem}[Closure of the perturbative weighted branch]\label{thm:weighted-closed}
For every \(m>\frac72\), the perturbative weighted ancient branch consists
only of the zero solution.
\end{theorem}

\begin{proof}
Let \(w\) be an ancient solution in the branch and fix \(\tau_0\in\mathbb R\).
For any \(s<\tau_0\), the time-shifted function
\[
   W_s(\theta)=w(s+\theta),\qquad \theta\ge0,
\]
is a global mild solution of \eqref{eq:gallay-vorticity} with initial datum
\(W_s(0)=w(s)\).  Since \(\|w(s)\|_m\le\rho_m\), \Cref{thm:gallay} applies to
this datum.  By uniqueness in
\(C([0,\infty);L^2_\sigma(m))\), this shifted ancient solution agrees with the
semiflow solution:
\[
   W_s(\theta)=\Phi_\theta w(s),\qquad \theta\ge0.
\]
Choosing \(\theta=\tau_0-s\) gives
\[
   w(\tau_0)=\Phi_{\tau_0-s}w(s).
\]
Applying \Cref{prop:weighted-lyap},
\[
   \mathcal L_m(w(\tau_0))
   \le
   e^{-2(\tau_0-s)}\mathcal L_m(w(s))
   \le
   C_m^2 e^{-2(\tau_0-s)}\|w(s)\|_m^2
   \le
   C_m^2\rho_m^2 e^{-2(\tau_0-s)} .
\]
The left-hand side is nonnegative and independent of \(s\).  As
\(s\to-\infty\), the factor \(e^{-2(\tau_0-s)}\) tends to zero.  Therefore
\(\mathcal L_m(w(\tau_0))=0\).  The coercive lower bound
\eqref{eq:weighted-coercive} gives
\[
   0\le \|w(\tau_0)\|_m^2\le \mathcal L_m(w(\tau_0))=0,
\]
so \(w(\tau_0)=0\) in \(L^2_\sigma(m)\).  Since \(\tau_0\) was arbitrary,
\(w\equiv0\) on \(\mathbb R\).
\end{proof}

\begin{remark}
This theorem closes only the perturbative weighted branch near zero vorticity.
It does not exclude a nonperturbative weighted ancient solution lying outside
the Gallay--Wayne small-data ball.
\end{remark}

\section{The Burgers-vortex model branch}

The next Lyapunov construction applies to the strained Burgers-vortex model.
It is a rigorous model theorem, not a theorem for the unforced equation
\eqref{eq:ns}.

\subsection{The strained vorticity equation}

Let \(x=(x_h,x_3)\in\mathbb R^2\times\mathbb R\) and let
\[
   A=\operatorname{diag}\left(-\frac12,-\frac12,1\right).
\]
The Burgers-vortex vorticity equation is
\begin{equation}\label{eq:burgers}
   \partial_\tau\Omega+(U\cdot\nabla)\Omega-(\Omega\cdot\nabla)U
   =
   \Delta\Omega-(Ax\cdot\nabla)\Omega+A\Omega,
   \qquad
   \nabla\cdot\Omega=0,
\end{equation}
with \(U\) recovered from \(\Omega\) by the Biot--Savart law.  It has the
stationary family
\begin{equation}\label{eq:burgers-family}
   \Omega=\alpha G,\qquad
   G(x)=\begin{pmatrix}0\\0\\g(x_h)\end{pmatrix},
   \qquad
   g(x_h)=\frac1{4\pi}e^{-|x_h|^2/4},
   \qquad \alpha\in\mathbb R .
\end{equation}

For \(m>1\), let \(L^2_h(m)\) denote the horizontal weighted space on
\(\mathbb R^2\) with norm
\[
   \|f\|_{L^2_h(m)}^2
   =
   \int_{\mathbb R^2}|f(x_h)|^2
   \left(1+\frac{|x_h|^2}{4m}\right)^m\,dx_h .
\]
Set
\[
   X(m)=BC(\mathbb R_{x_3};L^2_h(m)),\qquad
   X_0(m)=
   \left\{
      f\in X(m):
      \int_{\mathbb R^2}f(x_h,x_3)\,dx_h=0
      \text{ for every }x_3
   \right\},
\]
and
\[
   \mathcal X_{\mathrm{full}}(m)=X(m)^3,\qquad
   \mathcal X_{\mathrm{mod}}(m)=X(m)\times X(m)\times X_0(m).
\]

\begin{definition}[Admissible Burgers-vortex solutions]\label{def:admissible-bv}
In the Burgers-vortex section, a mild solution
\(\Omega:\mathbb R\to\mathcal X_{\mathrm{full}}(m)\) is called admissible if
the following operations are justified for every compact time interval:
\begin{enumerate}[label=\textup{(\roman*)}]
\item horizontal integrations by parts against constants and against
\(x_h\) are valid, with no boundary contribution at \(|x_h|=\infty\);
\item the identity
\[
   \int_{\mathbb R^2}
   \bigl[(U\cdot\nabla)\Omega-(\Omega\cdot\nabla)U\bigr]_3\,dx_h=0
\]
holds in distributions in \(x_3\) and \(\tau\);
\item the map
\[
   \tau\mapsto
   \int_{\mathbb R^2}\Omega_3(x_h,x_3,\tau)\,dx_h
\]
is absolutely continuous after testing in \(x_3\).
\end{enumerate}
Smooth rapidly decaying solutions are admissible.  The Gallay--Maekawa mild
solutions used below are assumed in this admissible class; this is precisely
the amount of regularity needed for the circulation argument.
\end{definition}

\begin{lemma}[Circulation coordinate]\label{lem:circulation}
Let \(m>1\) and let
\(\Omega\in\mathcal X_{\mathrm{full}}(m)\) be divergence-free.  Assume that
the horizontal integration by parts
\[
   \int_{\mathbb R^2}\nabla_h\cdot\Omega_h(x_h,x_3)\,dx_h=0
\]
is valid for every \(x_3\), or in the sense of distributions in \(x_3\).  Then
\[
   \alpha[\Omega](x_3)
   =
   \int_{\mathbb R^2}\Omega_3(x_h,x_3)\,dx_h
\]
is well-defined and independent of \(x_3\).  Writing the common value as
\(\alpha[\Omega]\), one has
\[
   Q\Omega:=\Omega-\alpha[\Omega]G\in\mathcal X_{\mathrm{mod}}(m).
\]
\end{lemma}

\begin{proof}
Since \(m>1\), the horizontal weighted space embeds into
\(L^1(\mathbb R^2)\).  Indeed, by Cauchy--Schwarz,
\[
   \int_{\mathbb R^2}|f(x_h)|\,dx_h
   \le
   \left(\int_{\mathbb R^2}|f(x_h)|^2
      \left(1+\frac{|x_h|^2}{4m}\right)^m dx_h\right)^{1/2}
   \left(\int_{\mathbb R^2}
      \left(1+\frac{|x_h|^2}{4m}\right)^{-m}dx_h\right)^{1/2}.
\]
The second integral is finite because the dimension is two and \(m>1\).  Hence
\(\alpha[\Omega](x_3)\) is well-defined.

Using \(\nabla\cdot\Omega=0\) gives
\[
   \partial_{x_3}\alpha[\Omega](x_3)
   =
   \int_{\mathbb R^2}\partial_{x_3}\Omega_3\,dx_h
   =
   -\int_{\mathbb R^2}\nabla_h\cdot\Omega_h\,dx_h
   =
   0,
\]
where the last equality is the stated horizontal integration-by-parts
assumption.  Thus \(\partial_{x_3}\alpha[\Omega]=0\), classically or
distributionally according to the interpretation of the previous display, and
\(\alpha[\Omega]\) is independent of \(x_3\).

Finally,
\[
   \int_{\mathbb R^2}g(x_h)\,dx_h
   =
   \frac1{4\pi}\int_0^\infty e^{-r^2/4}\,2\pi r\,dr
   =
   \frac12\int_0^\infty e^{-r^2/4}r\,dr
   =
   1 .
\]
Therefore the third component of \(Q\Omega=\Omega-\alpha[\Omega]G\) has zero
horizontal mean for every \(x_3\), while the first two components are
unchanged.  Hence \(Q\Omega\in\mathcal X_{\mathrm{mod}}(m)\).
\end{proof}

\begin{lemma}[Conservation of circulation]\label{lem:circulation-conserved}
Let \(m>1\), and let \(\Omega(\tau)\) be an admissible mild solution of
\eqref{eq:burgers} with values in \(\mathcal X_{\mathrm{full}}(m)\).  Then
\(\alpha[\Omega(\tau)]\) is independent of \(\tau\).  If
\(\bar\alpha\) denotes this constant, then
\[
   \omega(\tau)=\Omega(\tau)-\bar\alpha G
\]
solves the perturbation equation around the fixed Burgers vortex
\(\bar\alpha G\) and belongs to \(\mathcal X_{\mathrm{mod}}(m)\) for every
\(\tau\).
\end{lemma}

\begin{proof}
By admissibility, the following horizontal integrations can be read either
classically for smooth rapidly decaying solutions or distributionally for
admissible mild solutions.  Integrating the third component of
\eqref{eq:burgers} over \(x_h\), the nonlinear term does not contribute because
\[
   (U\cdot\nabla)\Omega-(\Omega\cdot\nabla)U
   =
   -\nabla\times(U\times\Omega),
\]
where the identity follows from
\[
   \nabla\times(U\times\Omega)
   =
   (\Omega\cdot\nabla)U-(U\cdot\nabla)\Omega
   +U(\nabla\cdot\Omega)-\Omega(\nabla\cdot U),
\]
and both \(U\) and \(\Omega\) are divergence-free.  The third component of a
curl is
\[
   [\nabla\times(U\times\Omega)]_3
   =
   \partial_{x_1}(U\times\Omega)_2
   -
   \partial_{x_2}(U\times\Omega)_1,
\]
a horizontal divergence.  Its integral over \(\mathbb R^2_{x_h}\) vanishes by
the horizontal decay.

For the linear part,
\[
   \bigl[\Delta\Omega-(Ax\cdot\nabla)\Omega+A\Omega\bigr]_3
   =
   \Delta_h\Omega_3+\partial_{x_3}^2\Omega_3
   +\frac12 x_h\cdot\nabla_h\Omega_3
   -x_3\partial_{x_3}\Omega_3+\Omega_3 .
\]
After horizontal integration, the \(\Delta_h\) term vanishes and
\[
   \int_{\mathbb R^2}\frac12 x_h\cdot\nabla_h\Omega_3\,dx_h
   =
   -\int_{\mathbb R^2}\Omega_3\,dx_h .
\]
Here the equality is obtained from
\[
   \nabla_h\cdot(x_h\Omega_3)
   =
   2\Omega_3+x_h\cdot\nabla_h\Omega_3
\]
and integration over \(\mathbb R^2\).  The last displayed term cancels with
the horizontal integral of \(+\Omega_3\).  Hence
\[
   \frac{d}{d\tau}\alpha[\Omega(\tau)]
   =
   \partial_{x_3}^2\alpha[\Omega(\tau)]
   -x_3\partial_{x_3}\alpha[\Omega(\tau)] .
\]
By \Cref{lem:circulation}, the right-hand side is zero.  Thus the circulation
is conserved.  Subtracting the stationary solution \(\bar\alpha G\) from
\eqref{eq:burgers} gives the fixed-circulation perturbation equation, and
\Cref{lem:circulation} gives membership in \(\mathcal X_{\mathrm{mod}}(m)\).
\end{proof}

\begin{theorem}[Gallay--Maekawa stability]\label{thm:gm}
Let \(m>2\) and \(\alpha\in\mathbb R\).  There exist constants
\(\delta_{\alpha,m}>0\) and \(C_{\alpha,m}\ge1\) such that every
divergence-free initial perturbation
\[
   \omega_0\in\mathcal X_{\mathrm{mod}}(m),
   \qquad
   \|\omega_0\|_{\mathcal X_{\mathrm{mod}}(m)}
   \le\delta_{\alpha,m},
\]
generates a unique global mild solution of the fixed-circulation perturbation
equation around \(\alpha G\), denoted
\(\Phi_\tau^\alpha\omega_0\), and
\[
   \|\Phi_\tau^\alpha\omega_0\|_{\mathcal X_{\mathrm{mod}}(m)}
   \le
   C_{\alpha,m}e^{-\tau/2}
   \|\omega_0\|_{\mathcal X_{\mathrm{mod}}(m)}
   \qquad(\tau\ge0).
\]
\end{theorem}

\begin{proof}
This is the Gallay--Maekawa nonlinear stability theorem for Burgers vortices in
the horizontal weighted space \cite{GallayMaekawa}.  It is used as an external
input.  The only conclusions used below are existence, uniqueness, the
fixed-circulation semiflow property, and the displayed exponential decay.
\end{proof}

\begin{proposition}[Profile-relative Lyapunov functional]\label{prop:bv-lyap}
Fix \(m>2\) and \(\alpha\in\mathbb R\).  On the ball
\[
   \|\omega_0\|_{\mathcal X_{\mathrm{mod}}(m)}
   \le\delta_{\alpha,m}
\]
define
\[
   \mathcal B_{\alpha,m}(\omega_0)
   =
   \sup_{\tau\ge0}e^\tau
   \|\Phi_\tau^\alpha\omega_0\|_{\mathcal X_{\mathrm{mod}}(m)}^2 .
\]
Then
\[
   \|\omega_0\|_{\mathcal X_{\mathrm{mod}}(m)}^2
   \le
   \mathcal B_{\alpha,m}(\omega_0)
   \le
   C_{\alpha,m}^2
   \|\omega_0\|_{\mathcal X_{\mathrm{mod}}(m)}^2,
\]
and
\[
   \mathcal B_{\alpha,m}(\Phi_s^\alpha\omega_0)
   \le
   e^{-s}\mathcal B_{\alpha,m}(\omega_0)
   \qquad(s\ge0).
\]
\end{proposition}

\begin{proof}
Gallay--Maekawa stability gives, for every \(\tau\ge0\),
\[
   e^\tau
   \|\Phi_\tau^\alpha\omega_0\|_{\mathcal X_{\mathrm{mod}}(m)}^2
   \le
   C_{\alpha,m}^2
   \|\omega_0\|_{\mathcal X_{\mathrm{mod}}(m)}^2 .
\]
Taking the supremum over \(\tau\ge0\) proves that
\(\mathcal B_{\alpha,m}\) is finite and gives the upper coercive bound.  The
lower bound follows from the value at \(\tau=0\):
\[
   \mathcal B_{\alpha,m}(\omega_0)
   \ge
   \|\Phi_0^\alpha\omega_0\|_{\mathcal X_{\mathrm{mod}}(m)}^2
   =
   \|\omega_0\|_{\mathcal X_{\mathrm{mod}}(m)}^2 .
\]

For \(s\ge0\), uniqueness of the fixed-circulation perturbation equation gives
the semiflow identity
\[
   \Phi_\tau^\alpha(\Phi_s^\alpha\omega_0)
   =
   \Phi_{\tau+s}^\alpha\omega_0 .
\]
Consequently,
\[
   \mathcal B_{\alpha,m}(\Phi_s^\alpha\omega_0)
   =
   \sup_{\tau\ge0}
      e^\tau
      \|\Phi_{\tau+s}^\alpha\omega_0\|_{\mathcal X_{\mathrm{mod}}(m)}^2 .
\]
Writing \(\sigma=\tau+s\), one gets
\[
   \mathcal B_{\alpha,m}(\Phi_s^\alpha\omega_0)
   =
   e^{-s}\sup_{\sigma\ge s}
      e^\sigma
      \|\Phi_\sigma^\alpha\omega_0\|_{\mathcal X_{\mathrm{mod}}(m)}^2
   \le
   e^{-s}\mathcal B_{\alpha,m}(\omega_0).
\]
\end{proof}

\begin{definition}[Burgers-vortex ancient tube]\label{def:bv-tube}
Fix \(m>2\).  The Burgers-vortex ancient tube consists of all admissible
ancient mild solutions \(\Omega:\mathbb R\to\mathcal X_{\mathrm{full}}(m)\) of
\eqref{eq:burgers} such that, with
\(\bar\alpha=\alpha[\Omega(\tau)]\),
\[
   \sup_{\tau\in\mathbb R}
   \|\Omega(\tau)-\bar\alpha G\|_{\mathcal X_{\mathrm{mod}}(m)}
   \le
   \delta_{\bar\alpha,m}.
\]
The phrase ``ancient mild'' also includes the compatibility that for every
\(s\in\mathbb R\), the shifted perturbation
\[
   \theta\mapsto \Omega(s+\theta)-\bar\alpha G,\qquad \theta\ge0,
\]
is the Gallay--Maekawa fixed-circulation mild solution launched from
\(\Omega(s)-\bar\alpha G\).
\end{definition}

\begin{theorem}[Ancient rigidity in the Burgers-vortex tube]\label{thm:bv-rigid}
Let \(m>2\).  If \(\Omega\) belongs to the Burgers-vortex ancient tube, then
\[
   \Omega(\tau)\equiv\bar\alpha G
   \qquad\text{for all }\tau\in\mathbb R,
\]
where \(\bar\alpha=\alpha[\Omega(\tau)]\).
\end{theorem}

\begin{proof}
By \Cref{lem:circulation-conserved}, \(\bar\alpha\) is independent of time and
\[
   \omega(\tau)=\Omega(\tau)-\bar\alpha G
\]
is an ancient solution of the fixed-circulation perturbation equation in the
small ball of \(\mathcal X_{\mathrm{mod}}(m)\).  Fix \(\tau_0\in\mathbb R\)
and \(s<\tau_0\).  The shifted function
\[
   \omega_s(\theta)=\omega(s+\theta),\qquad \theta\ge0,
\]
is a global mild solution of the fixed-circulation perturbation equation with
initial datum \(\omega_s(0)=\omega(s)\).  The tube hypothesis gives
\[
   \|\omega(s)\|_{\mathcal X_{\mathrm{mod}}(m)}
   \le \delta_{\bar\alpha,m},
\]
so Gallay--Maekawa uniqueness applies.  Therefore
\[
   \omega(\tau_0)=\Phi_{\tau_0-s}^{\bar\alpha}\omega(s).
\]
Using \Cref{prop:bv-lyap},
\[
   \mathcal B_{\bar\alpha,m}(\omega(\tau_0))
   \le
   e^{-(\tau_0-s)}
   \mathcal B_{\bar\alpha,m}(\omega(s))
   \le
   e^{-(\tau_0-s)}
   C_{\bar\alpha,m}^2
   \|\omega(s)\|_{\mathcal X_{\mathrm{mod}}(m)}^2 .
\]
The tube bound makes the last norm at most \(\delta_{\bar\alpha,m}^2\),
uniformly in \(s\).  Letting \(s\to-\infty\) gives
\(\mathcal B_{\bar\alpha,m}(\omega(\tau_0))=0\).  By the lower coercive bound
in \Cref{prop:bv-lyap},
\[
   0\le
   \|\omega(\tau_0)\|_{\mathcal X_{\mathrm{mod}}(m)}^2
   \le
   \mathcal B_{\bar\alpha,m}(\omega(\tau_0))
   =
   0 .
\]
Thus \(\omega(\tau_0)=0\).  Since \(\tau_0\) was arbitrary,
\(\Omega(\tau)=\bar\alpha G\) for every \(\tau\).
\end{proof}

\begin{assumption}[Reduction to the Burgers-vortex model]\label{ass:bv-reduction}
Let \(\mathcal C\) be a class of ancient solutions of an unforced or
renormalized Navier--Stokes problem.  Assume there are \(m>2\), a map
\[
   \mathcal T:\mathcal C\to
   C_{\mathrm{loc}}(\mathbb R;\mathcal X_{\mathrm{full}}(m)),
\]
and a nonnegative defect \(d[u(\tau)]\) such that for every \(u\in\mathcal C\):
\begin{enumerate}[label=\textup{(\roman*)}]
\item \(\Omega=\mathcal T[u]\) is an admissible ancient mild solution of
\eqref{eq:burgers} in the sense of \Cref{def:admissible-bv};
\item with \(\bar\alpha=\alpha[\Omega]\),
\[
   \|\Omega(\tau)-\bar\alpha G\|_{\mathcal X_{\mathrm{mod}}(m)}
   \le C_* d[u(\tau)]
   \qquad(\tau\in\mathbb R);
\]
\item \(\sup_{\tau\in\mathbb R}d[u(\tau)]\le
\delta_{\bar\alpha,m}/C_*\).
\end{enumerate}
\end{assumption}

\begin{corollary}[Conditional model import]\label{cor:bv-import}
Under \Cref{ass:bv-reduction}, every \(u\in\mathcal C\) is mapped by
\(\mathcal T\) to an exact Burgers vortex:
\[
   \mathcal T[u](\tau)\equiv\bar\alpha G .
\]
\end{corollary}

\begin{proof}
Let \(\Omega=\mathcal T[u]\).  By assumption (i), \(\Omega\) is an ancient mild
solution of \eqref{eq:burgers}.  By \Cref{lem:circulation-conserved}, its
circulation \(\bar\alpha=\alpha[\Omega]\) is constant in time.  Assumptions
(ii) and (iii) give, for every \(\tau\in\mathbb R\),
\[
   \|\Omega(\tau)-\bar\alpha G\|_{\mathcal X_{\mathrm{mod}}(m)}
   \le
   C_* d[u(\tau)]
   \le
   \delta_{\bar\alpha,m}.
\]
Thus \(\Omega\) belongs to the Burgers-vortex ancient tube.  Applying
\Cref{thm:bv-rigid} gives \(\Omega(\tau)\equiv\bar\alpha G\), which is the
claimed conclusion.
\end{proof}

\begin{remark}
\Cref{cor:bv-import} is intentionally conditional.  The strained equation
\eqref{eq:burgers} contains the imposed axial strain \(A\), whereas the
unforced Navier--Stokes equations do not.  A separate analytic reduction is
needed before Burgers-vortex rigidity can be used for an unforced singularity
core.
\end{remark}

\section{The scalar \texorpdfstring{\(\Gamma\)}{Gamma} defect}

We now turn to the remaining axisymmetric bounded-swirl branch.  The rigorous
output of this section is an exact scalar identity and a precise shell-control
criterion.  The criterion is not known to hold in full generality.

\subsection{The equation for \texorpdfstring{\(\Gamma\)}{Gamma}}

\begin{lemma}[Swirl equation]\label{lem:gamma-equation}
Let \((u,p)\) be a smooth axisymmetric solution of \eqref{eq:ns}.  Then
\(\Gamma=ru_\theta\) satisfies
\begin{equation}\label{eq:gamma}
   \partial_t\Gamma+b\cdot\nabla\Gamma+\frac2r\partial_r\Gamma
   =
   \Delta\Gamma,
   \qquad b=u_r e_r+u_z e_z .
\end{equation}
\end{lemma}

\begin{proof}
For an axisymmetric vector field, the vector Laplacian has angular component
\[
   (\Delta u)_\theta
   =
   \left(\partial_r^2+\frac1r\partial_r+\partial_z^2-\frac1{r^2}\right)
   u_\theta ,
\]
and the convective derivative has angular component
\[
   [(u\cdot\nabla)u]_\theta
   =
   u_r\partial_r u_\theta+u_z\partial_z u_\theta+\frac{u_r}{r}u_\theta .
\]
The pressure has no angular component because the flow is axisymmetric:
\((\nabla p)_\theta=r^{-1}\partial_\theta p=0\).  Therefore the angular
component of \eqref{eq:ns} is
\[
   \partial_tu_\theta+u_r\partial_r u_\theta+u_z\partial_z u_\theta
   +\frac{u_r}{r}u_\theta
   =
   \left(\partial_r^2+\frac1r\partial_r+\partial_z^2-\frac1{r^2}\right)
   u_\theta .
\]
Multiplying by \(r\) and using
\[
   r\partial_r u_\theta=\partial_r(ru_\theta)-u_\theta
   =\partial_r\Gamma-u_\theta
\]
gives cancellation of the zeroth-order transport term:
\[
   r u_r\partial_r u_\theta+u_r u_\theta
   =
   u_r\partial_r\Gamma .
\]
For the diffusion,
\[
   r\left(\partial_r^2+\frac1r\partial_r-\frac1{r^2}\right)u_\theta
   =
   \partial_r^2\Gamma-\frac1r\partial_r\Gamma .
\]
Indeed, since \(u_\theta=\Gamma/r\),
\[
   \partial_r u_\theta
   =
   \frac1r\partial_r\Gamma-\frac1{r^2}\Gamma
\]
and
\[
   \partial_r^2u_\theta
   =
   \frac1r\partial_r^2\Gamma
   -\frac2{r^2}\partial_r\Gamma
   +\frac2{r^3}\Gamma .
\]
Thus
\[
   r\partial_r^2u_\theta+\partial_r u_\theta-\frac1r u_\theta
   =
   \partial_r^2\Gamma-\frac1r\partial_r\Gamma .
\]
Since
\[
   \Delta\Gamma
   =
   \partial_r^2\Gamma+\frac1r\partial_r\Gamma+\partial_z^2\Gamma
\]
for axisymmetric scalar functions, the radial diffusion term is
\[
   \partial_r^2\Gamma-\frac1r\partial_r\Gamma+\partial_z^2\Gamma
   =
   \Delta\Gamma-\frac2r\partial_r\Gamma .
\]
Combining the displayed identities gives \eqref{eq:gamma}.
\end{proof}

\begin{proposition}[Weighted scalar defect identity]\label{prop:gamma-identity}
Let \(u\) be a smooth axisymmetric solution of \eqref{eq:ns}.  Set
\(W=\Gamma\), and let \(\phi=\phi(r,z,t)\ge0\) be a smooth compactly supported
axisymmetric weight.  Then
\begin{equation}\label{eq:gamma-identity}
   \frac12\frac{d}{dt}\int_{\mathbb R^3}W^2\phi\,dx
   +
   \int_{\mathbb R^3}|\nabla W|^2\phi\,dx
   =
   \frac12\int_{\mathbb R^3}W^2
   \left(
      \partial_t\phi+\Delta\phi+b\cdot\nabla\phi+\frac2r\partial_r\phi
   \right)\,dx .
\end{equation}
If \(W=\Gamma-c\), the same identity holds provided either \(c=0\) or
\(\phi\) vanishes in a neighbourhood of the axis.  Without this restriction,
the right-hand side in \eqref{eq:gamma-identity} must be augmented by the axis
boundary term
\[
   2\pi c^2\int_{\mathbb R}\phi(0,z,t)\,dz .
\]
\end{proposition}

\begin{proof}
Since \(W=\Gamma\), \Cref{lem:gamma-equation} gives
\[
   \partial_tW+b\cdot\nabla W+\frac2r\partial_rW=\Delta W .
\]
Equivalently,
\[
   \partial_tW
   =
   \Delta W-b\cdot\nabla W-\frac2r\partial_rW .
\]
Multiply by \(W\phi\) and integrate over \(\mathbb R^3\).  We compute each
term separately.  The time derivative gives
\[
   \int W\partial_tW\,\phi\,dx
   =
   \frac12\frac{d}{dt}\int W^2\phi\,dx
   -
   \frac12\int W^2\partial_t\phi\,dx .
\]
For the Laplacian,
\[
   \int W\Delta W\,\phi\,dx
   =
   -\int|\nabla W|^2\phi\,dx
   -\int W\nabla W\cdot\nabla\phi\,dx .
\]
Since \(W\nabla W=\frac12\nabla(W^2)\), compact support of \(\phi\) gives
\[
   -\int W\nabla W\cdot\nabla\phi\,dx
   =
   -\frac12\int\nabla(W^2)\cdot\nabla\phi\,dx
   =
   \frac12\int W^2\Delta\phi\,dx .
\]
Because \(u\) is divergence-free and \(b\) is the meridional part of an
axisymmetric divergence-free vector field,
\[
   \nabla\cdot b=0
\]
in the full three-dimensional divergence sense.  Explicitly,
\[
   \nabla\cdot b
   =
   \frac1r\partial_r(ru_r)+\partial_z u_z
   =
   \nabla\cdot u
   =
   0 .
\]
Hence
\[
   -\int W(b\cdot\nabla W)\phi\,dx
   =
   -\frac12\int b\cdot\nabla(W^2)\phi\,dx
   =
   \frac12\int W^2 b\cdot\nabla\phi\,dx .
\]
Finally,
\[
   -\int W\frac2r\partial_rW\,\phi\,dx
   =
   -\int \frac1r\partial_r(W^2)\phi\,dx .
\]
In cylindrical variables \(dx=r\,dr\,d\theta\,dz\), this term equals
\[
   -\int_0^{2\pi}\int_{\mathbb R}\int_0^\infty
      \partial_r(W^2)\phi\,dr\,dz\,d\theta .
\]
There is no boundary contribution at \(r=\infty\), because \(\phi\) is
compactly supported.  At the axis, one can first insert a smooth radial cutoff
which vanishes for \(r<\delta\), perform the integration by parts on
\(\delta<r<\infty\), and then let \(\delta\downarrow0\).  Smoothness of a
bounded axisymmetric velocity gives \(u_\theta(r,z,t)=O(r)\) near \(r=0\);
hence \(\Gamma(r,z,t)=ru_\theta(r,z,t)=O(r^2)\).  Therefore
\(W^2\phi=\Gamma^2\phi=O(r^4)\) at the axis, and the boundary term converges
to zero.  Thus
\[
   -\int_0^{2\pi}\int_{\mathbb R}\int_0^\infty
      \partial_r(W^2)\phi\,dr\,dz\,d\theta
   =
   \int_0^{2\pi}\int_{\mathbb R}\int_0^\infty
      W^2\partial_r\phi\,dr\,dz\,d\theta
   =
   \int W^2\frac1r\partial_r\phi\,dx .
\]
Therefore
\[
   -\int W\frac2r\partial_rW\,\phi\,dx
   =
   \int W^2\frac1r\partial_r\phi\,dx
   =
   \frac12\int W^2\frac2r\partial_r\phi\,dx .
\]
Putting the time, diffusion, transport, and singular radial terms together
gives \eqref{eq:gamma-identity}.

If \(W=\Gamma-c\), the differential equation is unchanged because derivatives
of \(c\) vanish.  The only altered step is the radial integration by parts at
the axis.  Since \(\Gamma(0,z,t)=0\), one has \(W(0,z,t)=-c\), and the missing
boundary contribution is
\[
   \int_0^{2\pi}\int_{\mathbb R}c^2\phi(0,z,t)\,dz\,d\theta
   =
   2\pi c^2\int_{\mathbb R}\phi(0,z,t)\,dz .
\]
This term vanishes if \(c=0\) or if \(\phi\) is supported away from the axis.
\end{proof}

\begin{remark}
For bounded ancient mild axisymmetric solutions with bounded \(\Gamma\), this
identity is not invoked automatically.  It is invoked only under
\Cref{ass:shell-control}, where smoothness, finiteness of the integrals, and
validity of the weighted identities for the chosen weights are assumed
explicitly.  No vortex-stretching term appears in \eqref{eq:gamma-identity};
the nonlinearity enters only through the drift \(b\cdot\nabla\phi\).
\end{remark}

\subsection{Anisotropic shell weights}

Let
\[
   \phi_{R,Z,\zeta}(r,z,t)=\lambda_R(r)\eta_Z(z-\zeta(t)),
\]
where \(\lambda_R,\eta_Z\) are smooth cutoffs and \(\zeta\) is a smooth axial
modulation.  Write \(\eta_Z=\eta_Z(z-\zeta(t))\) in formulas below.

\begin{proposition}[Shell-error decomposition]\label{prop:shell}
Let \(W=\Gamma\).  With the notation above,
\begin{equation}\label{eq:shell-identity}
   \frac12\frac{d}{dt}\int W^2\lambda_R\eta_Z\,dx
   +
   \int|\nabla W|^2\lambda_R\eta_Z\,dx
   =
   \frac12\int W^2\mathcal E_{R,Z,\zeta}[u]\,dx,
\end{equation}
where
\begin{equation}\label{eq:shell-error}
   \mathcal E_{R,Z,\zeta}[u]
   =
   \eta_Z\left(
      \lambda_R''+\frac3r\lambda_R'+u_r\lambda_R'
   \right)
   +
   \lambda_R\left(
      \eta_Z''+(u_z-\dot\zeta)\eta_Z'
   \right).
\end{equation}
\end{proposition}

\begin{proof}
For an axisymmetric scalar,
\[
   \Delta(\lambda_R\eta_Z)
   =
   \left(\lambda_R''+\frac1r\lambda_R'\right)\eta_Z
   +
   \lambda_R\eta_Z'' .
\]
This is the scalar cylindrical Laplacian
\[
   \Delta f=\partial_r^2f+\frac1r\partial_rf+\partial_z^2f
\]
for functions independent of \(\theta\), applied to the separated product
\(\lambda_R(r)\eta_Z(z-\zeta(t))\).  There are no mixed derivatives because
\(\lambda_R\) has no \(z\)-dependence and \(\eta_Z\) has no \(r\)-dependence.
Also
\[
   \frac2r\partial_r(\lambda_R\eta_Z)
   =
   \frac2r\lambda_R'\eta_Z ,
\]
and
\[
   b\cdot\nabla(\lambda_R\eta_Z)
   =
   u_r\lambda_R'\eta_Z+u_z\lambda_R\eta_Z' .
\]
Since the axial argument is \(z-\zeta(t)\),
\[
   \partial_t(\lambda_R\eta_Z)=-\dot\zeta\,\lambda_R\eta_Z' .
\]
Adding these four contributions gives
\[
   \partial_t\phi+\Delta\phi+b\cdot\nabla\phi+\frac2r\partial_r\phi
   =
   \eta_Z\left(\lambda_R''+\frac3r\lambda_R'+u_r\lambda_R'\right)
   +
   \lambda_R\left(\eta_Z''+(u_z-\dot\zeta)\eta_Z'\right).
\]
Substitution into \Cref{prop:gamma-identity} proves the formula.
\end{proof}

\begin{remark}[The radial geometry]
The differential expression
\[
   \lambda_R''+\frac3r\lambda_R'
\]
is the radial Laplacian in four dimensions.  Thus the natural radial barrier
for the \(\Gamma\)-defect is adapted to
\(\partial_r^2+3r^{-1}\partial_r\).  The shell errors in
\eqref{eq:shell-error} are precisely radial transport through
\(u_r\lambda_R'\), axial drift mismatch through
\((u_z-\dot\zeta)\eta_Z'\), and the two cutoff-curvature terms.
\end{remark}

\subsection{A coercive shell-control criterion}

The following statement is deliberately conditional.  It spells out a set of
quantitative hypotheses under which the scalar defect identity forces
\(\Gamma\) to vanish.  The difficult open step is proving these
hypotheses for the whole nonperiodic bounded-\(\Gamma\) ancient branch.

\begin{assumption}[Coercive shell control]\label{ass:shell-control}
Let \(u\) be a bounded ancient mild axisymmetric solution with bounded
\(\Gamma\).  Assume, in addition, that \(u\) is smooth on
\(\mathbb R^3\times(-\infty,0)\) and that the identities
\Cref{lem:gamma-equation,prop:gamma-identity,prop:shell} hold for the weights
below.  There exist expanding smooth weights
\[
   \phi_n(r,z,t)=\lambda_n(r)\eta_n(z-\zeta_n(t)),
   \qquad 0\le\phi_n\le1,\qquad \phi_n\uparrow1
\]
locally uniformly on compact subsets of
\(\mathbb R^3\times(-\infty,0)\), where \(\lambda_n,\eta_n,\zeta_n\) are
smooth and \(\phi_n\) is compactly supported in space for each fixed time.
There are constants \(a_n>0\), \(M_n<\infty\), \(\varepsilon_n\ge0\), with
\(\varepsilon_n/a_n\to0\), such that
\[
   L_n(t):=\int_{\mathbb R^3}\Gamma^2\phi_n(t)\,dx<\infty
\]
for every \(t<0\), \(L_n\) is locally absolutely continuous on
\((-\infty,0)\), and for every compact interval \([s,t]\subset(-\infty,0)\)
all integrals below are finite.  For all \(s<t<0\),
\begin{equation}\label{eq:shell-coercive}
   \int_{\mathbb R^3}|\nabla\Gamma|^2\phi_n(\sigma)\,dx
   \ge
   a_n L_n(\sigma)
   \qquad\text{for a.e. }\sigma\in(s,t),
\end{equation}
\begin{equation}\label{eq:shell-error-bound}
   \left|
   \int_{\mathbb R^3}\Gamma^2
   \mathcal E_n[u]\,dx
   \right|
   \le
   2\varepsilon_n
   \qquad\text{for a.e. }\sigma\in(s,t),
\end{equation}
and
\begin{equation}\label{eq:shell-past-bound}
   \sup_{\sigma<t}L_n(\sigma)\le M_n
   \qquad\text{for every fixed }t<0.
\end{equation}
Here \(\mathcal E_n[u]\) is the shell-error density
\eqref{eq:shell-error} associated with \(\phi_n\).
\end{assumption}

\begin{theorem}[Conditional scalar flattening]\label{thm:scalar-flattening}
Let \(u\) be a bounded ancient mild axisymmetric solution with bounded
\(\Gamma\).  If \Cref{ass:shell-control} holds, then
\[
   \Gamma(x,t)=0
   \qquad\text{for every }(x,t)\in\mathbb R^3\times(-\infty,0).
\]
\end{theorem}

\begin{proof}
Fix \(n\) and abbreviate \(W=\Gamma\).  Integrating
\eqref{eq:shell-identity} from \(s\) to \(t\), using
\eqref{eq:shell-coercive} and \eqref{eq:shell-error-bound}, gives
\[
   \frac12L_n(t)-\frac12L_n(s)
   +
   a_n\int_s^tL_n(\sigma)\,d\sigma
   \le
   \varepsilon_n(t-s).
\]
Indeed, the shell identity gives
\[
   \frac12 L_n(t)-\frac12L_n(s)
   +
   \int_s^t\int_{\mathbb R^3}|\nabla\Gamma|^2\phi_n\,dx\,d\sigma
   =
   \frac12\int_s^t\int_{\mathbb R^3}\Gamma^2\mathcal E_n[u]\,dx\,d\sigma .
\]
The coercivity assumption bounds the dissipation from below by
\[
   \int_s^t\int_{\mathbb R^3}|\nabla\Gamma|^2\phi_n\,dx\,d\sigma
   \ge
   a_n\int_s^tL_n(\sigma)\,d\sigma,
\]
and the error assumption bounds the right-hand side above by
\(\varepsilon_n(t-s)\).

Since \(L_n\) is locally absolutely continuous by \Cref{ass:shell-control}, the
preceding integral inequality is equivalent to the differential inequality, for
a.e. \(t<0\),
\[
   \frac{d}{dt}L_n(t)+2a_nL_n(t)\le2\varepsilon_n .
\]
To see this explicitly, apply the integral inequality on \((t,t+h)\), divide
by \(h\), and let \(h\downarrow0\) at a Lebesgue point of \(L_n'\) and
\(L_n\).  Multiplying by the integrating factor \(e^{2a_nt}\) gives, for a.e.
\(t\),
\[
   \frac{d}{dt}\left(e^{2a_nt}L_n(t)\right)
   =
   e^{2a_nt}\left(\frac{d}{dt}L_n(t)+2a_nL_n(t)\right)
   \le
   2\varepsilon_n e^{2a_nt}.
\]
Integrating this inequality from \(s\) to \(t\) yields
\[
   L_n(t)
   \le
   e^{-2a_n(t-s)}L_n(s)
   +
   \frac{\varepsilon_n}{a_n}
   \left(1-e^{-2a_n(t-s)}\right).
\]
By \eqref{eq:shell-past-bound}, \(L_n(s)\le M_n\) for \(s<t\).  Letting
\(s\to-\infty\) gives
\[
   L_n(t)\le\frac{\varepsilon_n}{a_n}.
\]
Now let \(n\to\infty\).  Since \(\varepsilon_n/a_n\to0\) and
\(\phi_n\uparrow1\) locally uniformly, for every compact set
\(K\subset\mathbb R^3\),
\[
   \int_K\Gamma(x,t)^2\,dx
   \le
   \liminf_{n\to\infty}L_n(t)
   =
   0
\]
for every \(t<0\) outside a null set.  The inequality follows because, for
each compact \(K\), local uniform convergence \(\phi_n\to1\) gives
\(\inf_K\phi_n(t)\to1\), and hence
\[
   \int_K\Gamma^2\,dx
   =
   \lim_{n\to\infty}\int_K\Gamma^2\phi_n\,dx
   \le
   \liminf_{n\to\infty}L_n(t).
\]
Thus \(\Gamma(\cdot,t)=0\) locally in \(L^2\) for a.e. \(t\).  By the
smoothness included in \Cref{ass:shell-control}, \(\Gamma\) is continuous on
\(\mathbb R^3\times(-\infty,0)\).  Let \(t_0<0\) be arbitrary, and choose
times \(t_j\to t_0\) outside the null set for which
\(\Gamma(\cdot,t_j)=0\) locally in \(L^2_x\).  For each \(j\), continuity in
space implies \(\Gamma(x,t_j)=0\) for every \(x\): otherwise nonzero value at
one point would give nonzero \(L^2\)-mass on a small ball.  Passing to the
limit \(j\to\infty\) and using continuity in time gives
\(\Gamma(x,t_0)=0\) for every \(x\in\mathbb R^3\).  Since \(t_0\) was
arbitrary, \(\Gamma\equiv0\) pointwise.
\end{proof}

\begin{corollary}[Reduction after scalar flattening]\label{cor:zero-gamma}
Let \(u\) be as in \Cref{thm:scalar-flattening}.  If in addition \(u\) belongs
to a class for which \(\Gamma\equiv0\) implies one of the Liouville
alternatives in \Cref{cor:axis-closed}, then \(u\) is Galilean-trivial.
\end{corollary}

\begin{proof}
\Cref{thm:scalar-flattening} gives \(\Gamma\equiv0\).  The additional
assumption places the solution in one of the already closed axisymmetric
branches.  The conclusion follows from \Cref{cor:axis-closed}.
\end{proof}

\begin{remark}
\Cref{thm:scalar-flattening} is not an unconditional closure of the
nonperiodic bounded-\(\Gamma\) branch.  It shows exactly what must be proved:
one needs shell weights for which the radial transport, axial drift mismatch,
and cutoff-curvature errors are small compared with the local coercivity
constant.  Without such control, the scalar identity remains a structural
identity rather than a Liouville theorem.
\end{remark}

\section{Summary of closed and open branches}

The rigorous exclusions obtained above are the following.

\begin{theorem}[Branch exclusion ledger]\label{thm:ledger}
The following ancient branches are closed.
\begin{enumerate}[label=\textup{(\roman*)}]
\item Bounded ancient axisymmetric unforced solutions satisfying any of the
four alternatives in \Cref{cor:axis-closed} are Galilean-trivial.
\item The perturbative weighted-vorticity branch of
\Cref{def:weighted-branch} is zero for every \(m>\frac72\).
\item The Burgers-vortex ancient tube of \Cref{def:bv-tube} consists exactly
of the stationary Burgers vortices \(\alpha G\).
\item Any bounded ancient axisymmetric solution satisfying the coercive
shell-control hypotheses of \Cref{ass:shell-control}, and for which
\(\Gamma\equiv0\) enters a closed axisymmetric Liouville class, is
Galilean-trivial.
\end{enumerate}
\end{theorem}

\begin{proof}
Item (i) is \Cref{cor:axis-closed}.  Item (ii) is
\Cref{thm:weighted-closed}.  Item (iii) is \Cref{thm:bv-rigid}.  Item (iv) is
\Cref{cor:zero-gamma}.
\end{proof}

\section{Analytic dependencies}

For clarity we record which assertions are proved in the paper and which are
imported as external PDE inputs.  This is part of the mathematical statement:
none of the imported assertions is used without being named, and none of the
conditional assumptions is converted into a theorem.

\begin{proposition}[Dependency ledger]\label{prop:dependency-ledger}
The conclusions of \Cref{thm:ledger} depend only on the following analytic
inputs.
\begin{enumerate}[label=\textup{(\roman*)}]
\item The axisymmetric exclusions in \Cref{thm:no-swirl,thm:pointwise,thm:lp-swirl,thm:periodic}
use exactly the corresponding Liouville theorems stated in
\Cref{thm:knss-input,thm:knss-pointwise-input,thm:lzz-input,thm:lrz-input},
together with the elementary Galilean reduction of
\Cref{lem:galilean}.
\item The perturbative weighted-vorticity exclusion uses the
Gallay--Wayne stability theorem stated in \Cref{thm:gallay}.  Apart from
that input, the proof uses only uniqueness of the weighted-vorticity semiflow,
the assumed compatibility of the ancient branch with that semiflow, and the
exponential decay estimate stated in \Cref{thm:gallay}.
\item The Burgers-vortex rigidity theorem uses the Gallay--Maekawa stability
theorem stated in \Cref{thm:gm}.  The reduction to a fixed-circulation
fiber uses the admissibility hypotheses of \Cref{def:admissible-bv}; in
particular, the horizontal integrations by parts and the conservation of
horizontal circulation are assumptions on the solution class unless they have
been verified by an independent regularity and decay argument.
\item The scalar flattening result uses no Liouville theorem before
\(\Gamma\) has vanished.  It uses the scalar identity of
\Cref{prop:shell}, the regularity and finiteness assertions built into
\Cref{ass:shell-control}, the coercive estimate
\eqref{eq:shell-coercive}, the shell-error estimate
\eqref{eq:shell-error-bound}, and the past bound
\eqref{eq:shell-past-bound}.
\item The final Galilean-trivial conclusion after scalar flattening requires
an additional inclusion of the resulting no-swirl solution in one of the
Liouville classes listed in \Cref{cor:axis-closed}.  The paper does not assert
that every bounded ancient axisymmetric solution with bounded \(\Gamma\)
satisfies the shell-control assumption or that every no-swirl solution belongs
to a specified external Liouville class without the hypotheses of that class.
\end{enumerate}
\end{proposition}

\begin{proof}
Each item follows by tracing the cited proof.  For (i), the proofs of
\Cref{thm:no-swirl,thm:pointwise,thm:lp-swirl,thm:periodic} invoke only the
corresponding external Liouville result before using
\Cref{lem:galilean} to remove constant vector solutions.  For (ii),
\Cref{thm:weighted-closed} first shifts the ancient solution to time
\(s\), applies the Gallay--Wayne semiflow from time \(s\) to time \(t\), and
then sends \(s\to-\infty\); the only analytic estimate used in that argument
is the exponential decay estimate stated in \Cref{thm:gallay}.  For (iii), \Cref{lem:circulation} and
\Cref{lem:circulation-conserved} use precisely the admissibility conditions
stated in \Cref{def:admissible-bv}, and \Cref{thm:bv-rigid} then applies
\Cref{thm:gm} on the corresponding fixed-circulation fiber.  For (iv),
\Cref{thm:scalar-flattening} integrates \eqref{eq:shell-identity}, applies
\eqref{eq:shell-coercive}, \eqref{eq:shell-error-bound}, and
\eqref{eq:shell-past-bound}, and passes to the increasing local limit
\(\phi_n\uparrow1\).  The final assertion is
\Cref{cor:zero-gamma}.  No further compactness, classification, or regularity
statement is used in these arguments.
\end{proof}

\begin{remark}[Remaining bounded-swirl obstruction]
The branch not closed by this paper is the full class of bounded ancient
axisymmetric unforced Navier--Stokes solutions with nontrivial bounded
\(\Gamma\), no \(L^p\) decay of \(\Gamma\), no axial periodicity, and no
verified shell-control estimate of the form \Cref{ass:shell-control}.  The
identity \eqref{eq:shell-identity} reduces this obstruction to explicit
transport terms, but does not eliminate them.
\end{remark}

\begin{thebibliography}{99}

\bibitem{CSTY2008}
C.-C. Chen, R. M. Strain, T.-P. Tsai, and H.-T. Yau,
\emph{Lower bound on the blow-up rate of the axisymmetric Navier--Stokes
equations},
Int. Math. Res. Not. IMRN \textbf{2008}, Art. ID rnn016, 31 pp.,
doi:10.1093/imrn/rnn016.

\bibitem{CSTY2009}
C.-C. Chen, R. M. Strain, T.-P. Tsai, and H.-T. Yau,
\emph{Lower bounds on the blow-up rate of the axisymmetric Navier--Stokes
equations II},
Comm. Partial Differential Equations \textbf{34} (2009), no. 3, 203--232,
doi:10.1080/03605300902793956.

\bibitem{ESS2003}
L. Escauriaza, G. A. Seregin, and V. Sverak,
\emph{\(L_{3,\infty}\)-solutions of the Navier--Stokes equations and backward
uniqueness},
Russian Math. Surveys \textbf{58} (2003), no. 2, 211--250,
doi:10.1070/RM2003v058n02ABEH000609.

\bibitem{GallayWayne2002}
T. Gallay and C. E. Wayne,
\emph{Long-time asymptotics of the Navier--Stokes and vorticity equations on
\(\mathbb R^3\)},
Philos. Trans. Roy. Soc. London Ser. A \textbf{360} (2002), no. 1799,
2155--2188, doi:10.1098/rsta.2002.1068.

\bibitem{GallayWayne2006}
T. Gallay and C. E. Wayne,
\emph{Three-dimensional stability of Burgers vortices: the low Reynolds number
case},
Physica D \textbf{213} (2006), no. 2, 164--180,
doi:10.1016/j.physd.2005.11.006.

\bibitem{GallayMaekawa}
T. Gallay and Y. Maekawa,
\emph{Three-dimensional stability of Burgers vortices},
Comm. Math. Phys. \textbf{302} (2011), no. 2, 477--511,
doi:10.1007/s00220-010-1132-6.

\bibitem{KNSS2009}
G. Koch, N. Nadirashvili, G. Seregin, and V. Sverak,
\emph{Liouville theorems for the Navier--Stokes equations and applications},
Acta Math. \textbf{203} (2009), no. 1, 83--105,
doi:10.1007/s11511-009-0039-6.

\bibitem{LeiZhang2011}
Z. Lei and Q. S. Zhang,
\emph{A Liouville theorem for the axially-symmetric Navier--Stokes equations},
J. Funct. Anal. \textbf{261} (2011), no. 8, 2323--2345,
doi:10.1016/j.jfa.2011.06.016.

\bibitem{LRZ}
Z. Lei, X. Ren, and Q. S. Zhang,
\emph{A Liouville theorem for axially symmetric Navier--Stokes equations on
\(\mathbb R^2\times\mathbb T^1\)},
Math. Ann. \textbf{383} (2022), no. 1--2, 415--431,
doi:10.1007/s00208-020-02128-9.

\bibitem{LZZ}
Z. Lei, Q. Zhang, and N. Zhao,
\emph{Improved Liouville theorems for axially symmetric Navier--Stokes
equations},
Sci. Sin. Math. \textbf{47} (2017), no. 10, 1183--1198,
doi:10.1360/N012016-00149.

\bibitem{SchmidRossi2004}
P. J. Schmid and M. Rossi,
\emph{Three-dimensional stability of a Burgers vortex},
J. Fluid Mech. \textbf{500} (2004), 103--112,
doi:10.1017/S0022112003007341.

\bibitem{SereginSverak2009}
G. Seregin and V. Sverak,
\emph{On Type I singularities of the local axi-symmetric solutions of the
Navier--Stokes equations},
Comm. Partial Differential Equations \textbf{34} (2009), no. 2, 171--201,
doi:10.1080/03605300802683687.

\end{thebibliography}

\end{document}
